\makeatletter
\def\input@path{{./}{../}{../../}{../../../}{../../../../}{preamble/}{../preamble/}{../../preamble/}{../../../preamble/}}
\makeatother
\input{preamble_exams.tex}
\pagestyle{empty}
\newcommand{\worklines}{\par\vspace{3mm}\hrule\vspace{4mm}\hrule\vspace{2mm}}
% Keep student problems intact; blanks provide two column-sized writing lines.
\newenvironment{studentproblem}{\par\addvspace{4pt}\noindent\begin{minipage}{\linewidth}}{\end{minipage}\par}
\newcommand{\filterblank}{\rule{0pt}{5mm}\rule{\linewidth}{0.4pt}}
\begin{document}
\newgeometry{left=\dimexpr2in/5\relax,right=\dimexpr2in/5\relax,top=\dimexpr7in/20\relax,bottom=\dimexpr7in/20\relax}
\begingroup\renewcommand{\baselinestretch}{1}\small
\setlength{\parskip}{2pt}\makeatletter\edef\arraystretch{\strip@pt\dimexpr105pt/100\relax}\makeatother
\begin{center}{\large\bfseries Grade 10 -- Term 1: C9 Classroom Activity}\end{center}



\noindent All payoffs and regrets are in thousands of dollars. EMV uses the stated summary values.
Apply each filter only to current survivors; keep summary values (including regret) fixed.
Maximax ranks Max highest first; Maximin ranks Min highest first; Minimax Regret ranks Regret lowest first; EMV/Bayes ranks EMV highest first.


\begin{studentproblem}
\textbf{1. City logistics policy.}\par
\textbf{(a) Complete the Filtering Table.} EMV: keep best 4; then Maximin: keep best 2; then Minimax Regret: select best.
\par\vspace{2pt}\noindent
\begingroup\setlength{\tabcolsep}{2.5pt}
\begin{minipage}[t]{0.30\linewidth}
\centering\textbf{Summary}\par\vspace{2pt}
\begin{tabular}{@{}lrrrr@{}}
\toprule
Alt. & EMV & Max & Min & Regret \\
\midrule
$B_1$ & 70 & 111 & 48 & 99 \\
$B_2$ & 62 & 66 & 54 & 96 \\
$B_3$ & 84 & 126 & 30 & 72 \\
$B_4$ & 86 & 150 & 30 & 111 \\
$B_5$ & 79 & 141 & 39 & 111 \\
$B_6$ & 90 & 102 & 66 & 84 \\
\bottomrule
\end{tabular}
\end{minipage}\hfill
\begin{minipage}[t]{0.68\linewidth}
\centering\textbf{Filtering}\par\vspace{2pt}
\begin{tabularx}{\linewidth}{@{}c >{\raggedright\arraybackslash}p{0.25\linewidth} >{\raggedright\arraybackslash}X >{\raggedright\arraybackslash}X@{}}
\toprule
Step & Rule & Survivors & Eliminated\\\midrule
0 & Initial & $B_1$, $B_2$, $B_3$, $B_4$, $B_5$, $B_6$ & ---\\
1 & EMV: best 4 & \filterblank & \filterblank\\
2 & Maximin: best 2 & \filterblank & \filterblank\\
3 & Regret: best 1 & \filterblank & \filterblank\\
\bottomrule\end{tabularx}
\end{minipage}
\endgroup\par\vspace{2pt}
\textbf{(b) Conceptual Understanding.} Why can $B_3$, whose Max Regret is 72, not reenter the final comparison?
\worklines
\end{studentproblem}

\begin{studentproblem}
\textbf{2. Regional distributor expansion.}\par
\textbf{(a) Complete the Filtering Table.} Maximax: keep best 3; then EMV: select best.
\par\vspace{2pt}\noindent
\begingroup\setlength{\tabcolsep}{2.5pt}
\begin{minipage}[t]{0.30\linewidth}
\centering\textbf{Summary}\par\vspace{2pt}
\begin{tabular}{@{}lrrrr@{}}
\toprule
Alt. & EMV & Max & Min & Regret \\
\midrule
$X_1$ & 26 & 42 & 18 & 108 \\
$X_2$ & 88 & 126 & 18 & 87 \\
$X_3$ & 85 & 147 & 48 & 66 \\
$X_4$ & 83 & 150 & 33 & 93 \\
$X_5$ & 76 & 120 & 45 & 105 \\
$X_6$ & 104 & 111 & 96 & 39 \\
\bottomrule
\end{tabular}
\end{minipage}\hfill
\begin{minipage}[t]{0.68\linewidth}
\centering\textbf{Filtering}\par\vspace{2pt}
\begin{tabularx}{\linewidth}{@{}c >{\raggedright\arraybackslash}p{0.25\linewidth} >{\raggedright\arraybackslash}X >{\raggedright\arraybackslash}X@{}}
\toprule
Step & Rule & Survivors & Eliminated\\\midrule
0 & Initial & $X_1$, $X_2$, $X_3$, $X_4$, $X_5$, $X_6$ & ---\\
1 & Maximax: best 3 & \filterblank & \filterblank\\
2 & EMV: best 1 & \filterblank & \filterblank\\
\bottomrule\end{tabularx}
\end{minipage}
\endgroup\par\vspace{2pt}
\textbf{(b) Conceptual Understanding.} Why is $X_6$, with the highest original EMV of 104, eliminated before EMV is applied?
\worklines
\end{studentproblem}

\begin{studentproblem}
\textbf{3. Trade finance guarantees.}\par
\textbf{(a) Complete the Filtering Table.} Maximin: keep best 5; then Minimax Regret: keep best 4; then EMV: keep best 2; then Maximax: select best.
\par\vspace{2pt}\noindent
\begingroup\setlength{\tabcolsep}{2.5pt}
\begin{minipage}[t]{0.30\linewidth}
\centering\textbf{Summary}\par\vspace{2pt}
\begin{tabular}{@{}lrrrr@{}}
\toprule
Alt. & EMV & Max & Min & Regret \\
\midrule
$G_1$ & 101 & 141 & 30 & 72 \\
$G_2$ & 94 & 150 & 39 & 108 \\
$G_3$ & 117 & 147 & 96 & 42 \\
$G_4$ & 70 & 78 & 66 & 81 \\
$G_5$ & 83 & 135 & 48 & 102 \\
$G_6$ & 52 & 102 & 21 & 126 \\
\bottomrule
\end{tabular}
\end{minipage}\hfill
\begin{minipage}[t]{0.68\linewidth}
\centering\textbf{Filtering}\par\vspace{2pt}
\begin{tabularx}{\linewidth}{@{}c >{\raggedright\arraybackslash}p{0.25\linewidth} >{\raggedright\arraybackslash}X >{\raggedright\arraybackslash}X@{}}
\toprule
Step & Rule & Survivors & Eliminated\\\midrule
0 & Initial & $G_1$, $G_2$, $G_3$, $G_4$, $G_5$, $G_6$ & ---\\
1 & Maximin: best 5 & \filterblank & \filterblank\\
2 & Regret: best 4 & \filterblank & \filterblank\\
3 & EMV: best 2 & \filterblank & \filterblank\\
4 & Maximax: best 1 & \filterblank & \filterblank\\
\bottomrule\end{tabularx}
\end{minipage}
\endgroup\par\vspace{2pt}
\textbf{(b) Conceptual Understanding.} Swap the first two criteria: Regret keeps 5, then Maximin keeps 4; leave the last two filters unchanged. Does the final choice change? Show the survivors.
\worklines
\end{studentproblem}
\newpage

\begin{studentproblem}
\textbf{4. Insurance portfolio mixes.}\par
\textbf{(a) Complete the Filtering Table.} Minimax Regret: keep best 4; then Maximax: keep best 2; then Maximin: select best.
\par\vspace{2pt}\noindent
\begingroup\setlength{\tabcolsep}{2.5pt}
\begin{minipage}[t]{0.30\linewidth}
\centering\textbf{Summary}\par\vspace{2pt}
\begin{tabular}{@{}lrrrr@{}}
\toprule
Alt. & EMV & Max & Min & Regret \\
\midrule
$Y_1$ & 51 & 69 & 21 & 129 \\
$Y_2$ & 99 & 135 & 30 & 120 \\
$Y_3$ & 117 & 150 & 63 & 72 \\
$Y_4$ & 124 & 150 & 75 & 60 \\
$Y_5$ & 122 & 135 & 114 & 36 \\
$Y_6$ & 68 & 111 & 24 & 111 \\
\bottomrule
\end{tabular}
\end{minipage}\hfill
\begin{minipage}[t]{0.68\linewidth}
\centering\textbf{Filtering}\par\vspace{2pt}
\begin{tabularx}{\linewidth}{@{}c >{\raggedright\arraybackslash}p{0.25\linewidth} >{\raggedright\arraybackslash}X >{\raggedright\arraybackslash}X@{}}
\toprule
Step & Rule & Survivors & Eliminated\\\midrule
0 & Initial & $Y_1$, $Y_2$, $Y_3$, $Y_4$, $Y_5$, $Y_6$ & ---\\
1 & Regret: best 4 & \filterblank & \filterblank\\
2 & Maximax: best 2 & \filterblank & \filterblank\\
3 & Maximin: best 1 & \filterblank & \filterblank\\
\bottomrule\end{tabularx}
\end{minipage}
\endgroup\par\vspace{2pt}
\textbf{(b) Conceptual Understanding.} If Step 2 kept the best 3 by Maximax instead of 2, would the final choice change? Use the new survivors to justify your answer.
\worklines
\end{studentproblem}

\begin{studentproblem}
\textbf{5. Solar project investment.}\par
\textbf{(a) Complete the Filtering Table.} Maximax: keep best 5; then Maximin: keep best 3; then Minimax Regret: select best.
\par\vspace{2pt}\noindent
\begingroup\setlength{\tabcolsep}{2.5pt}
\begin{minipage}[t]{0.30\linewidth}
\centering\textbf{Summary}\par\vspace{2pt}
\begin{tabular}{@{}lrrrr@{}}
\toprule
Alt. & EMV & Max & Min & Regret \\
\midrule
$S_1$ & 78 & 87 & 69 & 72 \\
$S_2$ & 112 & 132 & 81 & 57 \\
$S_3$ & 105 & 138 & 60 & 90 \\
$S_4$ & 77 & 102 & 63 & 72 \\
$S_5$ & 86 & 129 & 60 & 90 \\
$S_6$ & 116 & 150 & 63 & 75 \\
\bottomrule
\end{tabular}
\end{minipage}\hfill
\begin{minipage}[t]{0.68\linewidth}
\centering\textbf{Filtering}\par\vspace{2pt}
\begin{tabularx}{\linewidth}{@{}c >{\raggedright\arraybackslash}p{0.25\linewidth} >{\raggedright\arraybackslash}X >{\raggedright\arraybackslash}X@{}}
\toprule
Step & Rule & Survivors & Eliminated\\\midrule
0 & Initial & $S_1$, $S_2$, $S_3$, $S_4$, $S_5$, $S_6$ & ---\\
1 & Maximax: best 5 & \filterblank & \filterblank\\
2 & Maximin: best 3 & \filterblank & \filterblank\\
3 & Regret: best 1 & \filterblank & \filterblank\\
\bottomrule\end{tabularx}
\end{minipage}
\endgroup\par\vspace{2pt}
\textbf{(b) Conceptual Understanding.} Among the three alternatives entering Step 3, compare the choice by Minimax Regret with the choice by Maximax. What risk attitude does each choice represent?
\worklines
\end{studentproblem}

\begin{studentproblem}
\textbf{6. Retail inventory contracts.}\par
\textbf{(a) Complete the Filtering Table.} Maximin: keep best 4; then EMV: select best.
\par\vspace{2pt}\noindent
\begingroup\setlength{\tabcolsep}{2.5pt}
\begin{minipage}[t]{0.30\linewidth}
\centering\textbf{Summary}\par\vspace{2pt}
\begin{tabular}{@{}lrrrr@{}}
\toprule
Alt. & EMV & Max & Min & Regret \\
\midrule
$R_1$ & 85 & 93 & 72 & 54 \\
$R_2$ & 37 & 66 & 18 & 120 \\
$R_3$ & 96 & 138 & 69 & 78 \\
$R_4$ & 53 & 99 & 27 & 114 \\
$R_5$ & 100 & 138 & 30 & 69 \\
$R_6$ & 83 & 147 & 51 & 87 \\
\bottomrule
\end{tabular}
\end{minipage}\hfill
\begin{minipage}[t]{0.68\linewidth}
\centering\textbf{Filtering}\par\vspace{2pt}
\begin{tabularx}{\linewidth}{@{}c >{\raggedright\arraybackslash}p{0.25\linewidth} >{\raggedright\arraybackslash}X >{\raggedright\arraybackslash}X@{}}
\toprule
Step & Rule & Survivors & Eliminated\\\midrule
0 & Initial & $R_1$, $R_2$, $R_3$, $R_4$, $R_5$, $R_6$ & ---\\
1 & Maximin: best 4 & \filterblank & \filterblank\\
2 & EMV: best 1 & \filterblank & \filterblank\\
\bottomrule\end{tabularx}
\end{minipage}
\endgroup\par\vspace{2pt}
\textbf{(b) Conceptual Understanding.} If the final EMV filter were replaced by Minimax Regret, would the selected alternative change? Compare the relevant values.
\worklines
\end{studentproblem}

\endgroup
% Restoring geometry clears the worksheet page and starts solutions.
\restoregeometry
\renewcommand{\baselinestretch}{1}\small
\setlength{\tabcolsep}{2.5pt}\setlength{\parskip}{2pt}
\section*{Solutions}
Parentheses show the statistic used at each filtering step; each elimination is new.
All summary values remain fixed. No cutoff or final ties occur.


\begin{studentproblem}
\textbf{1. City logistics policy}\par
\textbf{(a)}
\begin{center}
\small
\begin{tabularx}{\textwidth}{@{}c >{\raggedright\arraybackslash}X >{\raggedright\arraybackslash}p{0.29\textwidth} >{\raggedright\arraybackslash}p{0.23\textwidth}@{}}
\toprule
Step & Rule & Survivors & Eliminated \\
\midrule
0 & Initial & $\{B_1,B_2,B_3,B_4,B_5,B_6\}$ & --- \\
1 & EMV: best 4 & $B_6(90)$, $B_4(86)$, $B_3(84)$, $B_5(79)$ & $B_1(70)$, $B_2(62)$ \\
2 & Maximin: best 2 & $B_6(66)$, $B_5(39)$ & $B_4(30)$, $B_3(30)$ \\
3 & Regret: best 1 & \textbf{Selected:} $B_6(84)$ & $B_5(111)$ \\
\bottomrule
\end{tabularx}
\end{center}
\textbf{(b)} $B_3$ is eliminated at Step 2: its minimum 30 is below the cutoff 39. Only $B_6$ and $B_5$ reach Step 3. Although $72<84$, eliminated alternatives cannot reenter.
\end{studentproblem}

\begin{studentproblem}
\textbf{2. Regional distributor expansion}\par
\textbf{(a)}
\begin{center}
\small
\begin{tabularx}{\textwidth}{@{}c >{\raggedright\arraybackslash}X >{\raggedright\arraybackslash}p{0.29\textwidth} >{\raggedright\arraybackslash}p{0.23\textwidth}@{}}
\toprule
Step & Rule & Survivors & Eliminated \\
\midrule
0 & Initial & $\{X_1,X_2,X_3,X_4,X_5,X_6\}$ & --- \\
1 & Maximax: best 3 & $X_4(150)$, $X_3(147)$, $X_2(126)$ & $X_5(120)$, $X_6(111)$, $X_1(42)$ \\
2 & EMV: best 1 & \textbf{Selected:} $X_2(88)$ & $X_3(85)$, $X_4(83)$ \\
\bottomrule
\end{tabularx}
\end{center}
\textbf{(b)} $X_6$ has maximum 111, below the Step 1 cutoff 126. The surviving EMVs are 83, 85, and 88, so $X_2$ wins with 88 despite $X_6$ having EMV 104.
\end{studentproblem}

\begin{studentproblem}
\textbf{3. Trade finance guarantees}\par
\textbf{(a)}
\begin{center}
\small
\begin{tabularx}{\textwidth}{@{}c >{\raggedright\arraybackslash}X >{\raggedright\arraybackslash}p{0.29\textwidth} >{\raggedright\arraybackslash}p{0.23\textwidth}@{}}
\toprule
Step & Rule & Survivors & Eliminated \\
\midrule
0 & Initial & $\{G_1,G_2,G_3,G_4,G_5,G_6\}$ & --- \\
1 & Maximin: best 5 & $G_3(96)$, $G_4(66)$, $G_5(48)$, $G_2(39)$, $G_1(30)$ & $G_6(21)$ \\
2 & Regret: best 4 & $G_3(42)$, $G_1(72)$, $G_4(81)$, $G_5(102)$ & $G_2(108)$ \\
3 & EMV: best 2 & $G_3(117)$, $G_1(101)$ & $G_5(83)$, $G_4(70)$ \\
4 & Maximax: best 1 & \textbf{Selected:} $G_3(147)$ & $G_1(141)$ \\
\bottomrule
\end{tabularx}
\end{center}
\textbf{(b)} Yes. The swapped filters retain $G_3,G_1,G_4,G_5,G_2$ by regret, then $G_3,G_4,G_5,G_2$ by minimum. EMV retains $G_3(117),G_2(94)$; Maximax selects $G_2(150)$ over $G_3(147)$. The original choice is $G_3$.
\end{studentproblem}
\newpage

\begin{studentproblem}
\textbf{4. Insurance portfolio mixes}\par
\textbf{(a)}
\begin{center}
\small
\begin{tabularx}{\textwidth}{@{}c >{\raggedright\arraybackslash}X >{\raggedright\arraybackslash}p{0.29\textwidth} >{\raggedright\arraybackslash}p{0.23\textwidth}@{}}
\toprule
Step & Rule & Survivors & Eliminated \\
\midrule
0 & Initial & $\{Y_1,Y_2,Y_3,Y_4,Y_5,Y_6\}$ & --- \\
1 & Regret: best 4 & $Y_5(36)$, $Y_4(60)$, $Y_3(72)$, $Y_6(111)$ & $Y_2(120)$, $Y_1(129)$ \\
2 & Maximax: best 2 & $Y_4(150)$, $Y_3(150)$ & $Y_5(135)$, $Y_6(111)$ \\
3 & Maximin: best 1 & \textbf{Selected:} $Y_4(75)$ & $Y_3(63)$ \\
\bottomrule
\end{tabularx}
\end{center}
\textbf{(b)} Yes. Step 2 would retain $Y_4(150),Y_3(150),Y_5(135)$. Maximin would select $Y_5(114)$ over $Y_4(75)$ and $Y_3(63)$, replacing the original choice $Y_4$. The tie at 150 is inside the retained set, not at a cutoff.
\end{studentproblem}

\begin{studentproblem}
\textbf{5. Solar project investment}\par
\textbf{(a)}
\begin{center}
\small
\begin{tabularx}{\textwidth}{@{}c >{\raggedright\arraybackslash}X >{\raggedright\arraybackslash}p{0.29\textwidth} >{\raggedright\arraybackslash}p{0.23\textwidth}@{}}
\toprule
Step & Rule & Survivors & Eliminated \\
\midrule
0 & Initial & $\{S_1,S_2,S_3,S_4,S_5,S_6\}$ & --- \\
1 & Maximax: best 5 & $S_6(150)$, $S_3(138)$, $S_2(132)$, $S_5(129)$, $S_4(102)$ & $S_1(87)$ \\
2 & Maximin: best 3 & $S_2(81)$, $S_6(63)$, $S_4(63)$ & $S_3(60)$, $S_5(60)$ \\
3 & Regret: best 1 & \textbf{Selected:} $S_2(57)$ & $S_4(72)$, $S_6(75)$ \\
\bottomrule
\end{tabularx}
\end{center}
\textbf{(b)} Minimax Regret selects $S_2(57)$ over $S_4(72),S_6(75)$, limiting the worst missed opportunity. Maximax selects $S_6(150)$ over $S_2(132),S_4(102)$, emphasizing the largest possible payoff. These risk attitudes lead to different choices.
\end{studentproblem}

\begin{studentproblem}
\textbf{6. Retail inventory contracts}\par
\textbf{(a)}
\begin{center}
\small
\begin{tabularx}{\textwidth}{@{}c >{\raggedright\arraybackslash}X >{\raggedright\arraybackslash}p{0.29\textwidth} >{\raggedright\arraybackslash}p{0.23\textwidth}@{}}
\toprule
Step & Rule & Survivors & Eliminated \\
\midrule
0 & Initial & $\{R_1,R_2,R_3,R_4,R_5,R_6\}$ & --- \\
1 & Maximin: best 4 & $R_1(72)$, $R_3(69)$, $R_6(51)$, $R_5(30)$ & $R_4(27)$, $R_2(18)$ \\
2 & EMV: best 1 & \textbf{Selected:} $R_5(100)$ & $R_3(96)$, $R_1(85)$, $R_6(83)$ \\
\bottomrule
\end{tabularx}
\end{center}
\textbf{(b)} Yes. Among $R_1,R_3,R_6,R_5$, EMV selects $R_5(100)$ over 85, 96, and 83. Regret instead selects $R_1(54)$ over $R_5(69),R_3(78),R_6(87)$.
\end{studentproblem}
\end{document}
